% =====================================================================
\chapter{Optimización bayesiana}
\label{tema:4}
% =====================================================================

La optimización bayesiana usa el sustituto probabilístico del \cref{tema:3} para decidir el siguiente experimento: elige el punto que maximiza una \emph{función de adquisición}, una medida del valor esperado de ensayarlo según el modelo. El método procede de Jones, Schonlau y Welch~\cite{jones}, que introdujeron el algoritmo EGO con la mejora esperada; la exposición sigue la §4 del tutorial de Frazier~\cite{frazier}, los capítulos~7, 8 y~11 de Garnett~\cite{garnett}, la revisión de Shahriari \emph{et al.}~\cite{shahriari} y el capítulo~16 de Kochenderfer y Wheeler~\cite{kw}. Todo el tema trata problemas de \emph{minimización}; para maximizar basta cambiar el signo de $f$.

\section{El algoritmo}

\begin{algoritmo}[Optimización bayesiana]\label{alg:bo}
Datos: dominio $\X$, presupuesto $N$, diseño inicial de $n_0$ puntos.
\begin{enumerate}[label=\arabic*.]
  \item Evaluar el diseño inicial (por ejemplo, un LHS con $n_0\approx 10\,d$ o menos si el presupuesto es corto) y formar $\D_{n_0}$.
  \item Para $n=n_0,\dots,N-1$:
  \begin{enumerate}[label=(\alph*)]
    \item ajustar el sustituto a $\D_n$ (hiperparámetros por verosimilitud marginal);
    \item calcular $x_{n+1}\in\arg\max_{x\in\X}\alpha_n(x)$ con un optimizador auxiliar (varios arranques de L-BFGS-B o una nube de candidatos);
    \item evaluar $y_{n+1}=f(x_{n+1})+\varepsilon$ y añadir el par a los datos.
  \end{enumerate}
  \item Recomendar $\hat x_N\in\arg\min_{x}\mu_N(x)$ (entre los puntos evaluados o en todo $\X$) y confirmar con ensayos repetidos.
\end{enumerate}
\end{algoritmo}

La función de adquisición debe equilibrar dos tendencias: \emph{explotar} (ensayar donde la media $\mu_n$ es baja) y \emph{explorar} (ensayar donde la incertidumbre $\sigma_n$ es grande). En lo que sigue, $\mu=\mu_n(x)$, $\sigma=\sigma_n(x)$ es la desviación típica de la función latente $f(x)$ (no de una nueva observación) y $\fmejor$ es el valor de referencia: el mínimo observado si no hay ruido o el mínimo de $\mu_n$ en los puntos evaluados si lo hay~\cite[§5]{frazier}.

\section{Probabilidad de mejora}

\begin{definicion}[Probabilidad de mejora]\label{def:pi}
$\PI_n(x)=\Pr\bigl[f(x)<\fmejor-\xi\bigr]=\Phi\Bigl(\dfrac{\fmejor-\xi-\mu}{\sigma}\Bigr)$, con un margen $\xi\ge 0$.
\end{definicion}

La igualdad se debe a que $f(x)\sim\N(\mu,\sigma^2)$ y a la tipificación (módulo~0.8, §6.3). La PI ignora \emph{cuánto} se mejora: prefiere una mejora casi segura y diminuta a una probable y grande, y con $\xi=0$ tiende a muestrear junto al mejor punto actual~\cite[§7.5]{garnett}.

\section{Mejora esperada}

\begin{definicion}[Mejora y mejora esperada]\label{def:ei}
La \emph{mejora} en $x$ es $I(x)=\max\bigl(\fmejor-f(x),0\bigr)$, y la \emph{mejora esperada} es $\EI_n(x)=\E_n[I(x)]$, donde la esperanza se toma respecto a la posterior de $f(x)$ dados los datos~\cite[ec.~(14)]{jones}.
\end{definicion}

\begin{teorema}[Forma cerrada de la mejora esperada]\label{teo:ei}
Si $f(x)\mid\D_n\sim\N(\mu,\sigma^2)$ con $\sigma>0$ y $z=(\fmejor-\mu)/\sigma$, entonces
\[
\EI_n(x)=(\fmejor-\mu)\,\Phi(z)+\sigma\,\varphi(z).
\]
Si $\sigma=0$, $\EI_n(x)=\max(\fmejor-\mu,0)$.
\end{teorema}

Frazier~\cite[ec.~(8)]{frazier} da una forma equivalente para maximización, escrita con $\Delta=\mu-\fmejor$ y $[\Delta]^+$; Jones \emph{et al.}~\cite[ec.~(15)]{jones} la dan en la forma del teorema.

\begin{proposicion}[Propiedades de la mejora esperada]\label{prop:ei}
Como función de $(\mu,\sigma)$, con $\sigma>0$:
\begin{enumerate}
  \item $\EI>0$ y $\EI\ge\max(\fmejor-\mu,0)$;
  \item $\partial\EI/\partial\mu=-\Phi(z)<0$: a igual incertidumbre, se prefiere la media más baja;
  \item $\partial\EI/\partial\sigma=\varphi(z)>0$: a igual media, se prefiere la mayor incertidumbre.
\end{enumerate}
\end{proposicion}

\addtocounter{definicion}{1}%

\addtocounter{figure}{1}%

\section{Cotas de confianza y muestreo de Thompson}

\begin{definicion}[Cota inferior de confianza]\label{def:lcb}
Para un parámetro $\beta_n>0$, se elige $x_{n+1}\in\arg\min_x\LCB_n(x)$, con $\LCB_n(x)=\mu_n(x)-\sqrt{\beta_n}\,\sigma_n(x)$. En maximización se usa la cota superior $\UCB_n=\mu_n+\sqrt{\beta_n}\sigma_n$ (GP-UCB).
\end{definicion}

Srinivas \emph{et al.}~\cite{srinivas} demostraron que, con una sucesión $\beta_n$ que crece logarítmicamente, el regret acumulado de GP-UCB crece como $O^*\bigl(\sqrt{N\beta_N\gamma_N}\bigr)$, donde $\gamma_N$ es la máxima información que $N$ observaciones ruidosas pueden aportar sobre $f$; para el núcleo exponencial cuadrático, $\gamma_N$ crece como una potencia de $\log N$, y el regret medio $R_N/N$ tiende a cero. La teoría da una justificación al equilibrio entre explorar y explotar, pero los valores de $\beta_n$ que exige son conservadores; en la práctica se usan constantes entre $1$ y $9$~\cite[§8.4 y cap.~10]{garnett}.

El \emph{muestreo de Thompson} extrae una función $\tilde f\sim\GP(\mu_n,k_n)$ de la posterior y ensaya su mínimo; explora en proporción a la probabilidad de que cada punto sea el óptimo y se paraleliza con facilidad~\cite[§7.9]{garnett}.

\section{Optimizar la función de adquisición}

La función $\alpha_n$ es barata pero multimodal: vale casi cero en la mayor parte del dominio y tiene picos estrechos. Lo habitual es evaluarla en una nube de miles de candidatos (aleatorios y perturbaciones de los mejores puntos) y refinar los mejores con L-BFGS-B, usando su gradiente analítico o por diferenciación automática~\cite[§9.2]{garnett}, \cite{botorch}. Con valores de EI muy pequeños conviene trabajar con su logaritmo para evitar regiones numéricamente planas.

\begin{observacion}[Varianza latente y varianza predictiva]\label{obs:latente}
Muchas bibliotecas devuelven la desviación típica de una nueva \emph{observación}, $\sqrt{\sigma_n^2(x)+\sigma_\varepsilon^2}$, en lugar de la de $f(x)$. Usarla en la EI o en la probabilidad de cumplir una restricción es un error: en un punto ya ensayado la incertidumbre nunca baja del nivel de ruido, la EI no se anula y el algoritmo repite indefinidamente el mismo ajuste. En la práctica, con \texttt{GaussianProcessRegressor} y un \texttt{WhiteKernel}, hay que restar la varianza del ruido. En los experimentos de la práctica con la inyectora, sin esta corrección el algoritmo dedicaba la segunda mitad del presupuesto a repetir ajustes casi idénticos en una esquina del dominio.
\end{observacion}

\section{Restricciones}

\begin{proposicion}[Mejora esperada con restricciones]\label{prop:eic}
Sea el problema de la \cref{def:restricciones} con una restricción $c(x)\le 0$ modelada con un proceso gaussiano independiente del de $f$, con posterior $\N(\mu_c,\sigma_c^2)$, y sea $\fmejor$ el mejor valor \emph{factible}. La mejora restringida $I_c(x)=I(x)\,\mathbf 1[c(x)\le 0]$ tiene esperanza
\[
\EI_c(x)=\EI_n(x)\cdot\Phi\Bigl(\frac{-\mu_c(x)}{\sigma_c(x)}\Bigr).
\]
\end{proposicion}

Gardner \emph{et al.}~\cite{gardner} y Gelbart, Snoek y Adams~\cite{gelbart} estudian la optimización bayesiana con restricciones desconocidas colocando también un proceso gaussiano sobre cada restricción; el segundo admite además restricciones que se evalúan por separado del objetivo. El producto de la \cref{prop:eic} es la forma más sencilla de combinar ambos modelos~\cite[§11.2]{garnett}. Para una especificación $w(x)\le w_{\max}$, la probabilidad de cumplirla es $\Phi\bigl((w_{\max}-\mu_w)/\sigma_w\bigr)$. Si todavía no hay ningún punto factible, se maximiza solo la probabilidad de factibilidad~\cite[§11.2]{garnett}. Las restricciones conocidas y baratas no se modelan: se imponen al generar los candidatos.

\section{Evaluaciones por lotes}

En una planta se suelen hacer varios ensayos por turno, o hay varias máquinas: se necesita proponer $q$ puntos a la vez. La generalización natural es la \emph{mejora esperada multipunto}~\cite{ginsbourger},
\[
q\text{-}\EI(x^{(1)},\dots,x^{(q)})=\E_n\Bigl[\max\bigl(\fmejor-\min_{j}f(x^{(j)}),0\bigr)\Bigr],
\]
que no tiene forma cerrada útil y se estima por Monte Carlo con la distribución conjunta de los $q$ valores; BoTorch la optimiza con gradientes de Monte Carlo~\cite{botorch}. Dos heurísticas secuenciales más sencillas, de Ginsbourger, Le Riche y Carraro~\cite{ginsbourger}, eligen los puntos uno a uno añadiendo un dato ficticio en cada punto elegido: el \emph{creyente del kriging} usa como valor la media posterior $\mu_n(x^{(j)})$, y el \emph{mentiroso constante} usa un valor fijo $L$ (el mínimo, la media o el máximo observado; cuanto mayor es $L$, más se dispersan los puntos). Como la varianza posterior no depende de los valores (\cref{obs:posterior}), el dato ficticio reduce la incertidumbre alrededor de $x^{(j)}$ y aleja de él al siguiente punto del lote. Si los ensayos terminan en momentos distintos, los pendientes se tratan igual (optimización asíncrona~\cite[§11.4]{garnett}).

\section{Varios objetivos}

En la inyectora se puede querer, en lugar de un único coste, minimizar a la vez la fracción defectuosa y el tiempo de ciclo, y dejar la elección final a la dirección de la planta.

\begin{definicion}[Dominancia y frente de Pareto]\label{def:pareto}
Sean $F(x)=(F_1(x),\dots,F_m(x))$ objetivos a minimizar. El punto $x$ \emph{domina} a $x'$ si $F_i(x)\le F_i(x')$ para todo $i$ y $F_i(x)<F_i(x')$ para algún $i$. Un punto es \emph{óptimo de Pareto} si ningún otro lo domina; la imagen del conjunto de óptimos es el \emph{frente de Pareto}~\cite[§12.1]{kw}, \cite[§9.2]{mn}.
\end{definicion}

\begin{definicion}[Hipervolumen]\label{def:hv}
Dado un punto de referencia $r$ dominado por todos los puntos de interés, el \emph{hipervolumen} de un conjunto $P$ de vectores objetivo es la medida (área en dos dimensiones) de la región $\{v: \exists p\in P,\ p\le v\le r\}$ dominada por $P$ y acotada por $r$.
\end{definicion}

El hipervolumen crece si y solo si se añade un punto no dominado por el conjunto, y por eso es la medida habitual de progreso. La \emph{mejora esperada del hipervolumen} (EHVI) es la generalización de la EI; Daulton, Balandat y Bakshy~\cite{qehvi} la extienden a lotes con restricciones (qEHVI) y la calculan con diferenciación automática. Una alternativa más sencilla es escalarizar.

\begin{proposicion}[Escalarización de Chebyshev]\label{prop:chebyshev}
Sean $w_i>0$ pesos y $z^\star$ un punto ideal con $z^\star_i\le F_i(x)$ para todo $x$. Todo minimizador de $s_w(x)=\max_i w_i\bigl(F_i(x)-z^\star_i\bigr)$ es débilmente óptimo de Pareto: no existe $x'$ con $F_i(x')<F_i(x)$ para todo $i$.
\end{proposicion}

ParEGO~\cite{parego} aplica la EI a la escalarización de Chebyshev aumentada con pesos aleatorios distintos en cada iteración, y así recorre el frente. A diferencia de la suma ponderada, la escalarización de Chebyshev alcanza también los puntos de las partes no convexas del frente (\cref{eje:4-chebyshev}).

\section{Varias fidelidades}

A menudo hay versiones más baratas e imprecisas del experimento: una simulación del llenado del molde antes del ensayo en máquina, un ensayo de resistencia a 7 días en lugar de 28, un entrenamiento corto en lugar de uno completo. La optimización multifidelidad modela conjuntamente $f(x,s)$, con $s$ la fidelidad, y elige a la vez el punto y la fidelidad~\cite[§11.5]{garnett}, \cite[§5]{frazier}. La regla más simple es dividir la adquisición por el coste: Snoek, Larochelle y Adams~\cite{snoek} proponen la mejora esperada por segundo, $\EI(x)/c(x)$, con el coste $c(x)$ también modelado. Kandasamy \emph{et al.}~\cite{kandasamy} (MF-GP-UCB) usan las fidelidades baratas para descartar regiones poco prometedoras y reservan las caras para las regiones buenas, con garantías de regret mejores que las de ignorar las aproximaciones.

\section{Caso de estudio: el optimizador de la inyectora}

\addtocounter{definicion}{1}%

\begin{table}[ht]
\centering\small
\begin{tabularx}{\textwidth}{@{}>{\raggedright\arraybackslash}p{2.5cm}>{\raggedright\arraybackslash}X>{\raggedright\arraybackslash}p{1.4cm}>{\raggedright\arraybackslash}p{3.8cm}@{}}
\toprule
Herramienta & Para qué sirve & Licencia & Cuándo elegirla \\
\midrule
BoTorch~\cite{botorch} & Adquisiciones de Monte Carlo: qEI, qNEI, qEHVI, multifidelidad, restricciones & MIT & Problemas a medida, lotes, varios objetivos \\
Ax~\cite{ax} & Interfaz de alto nivel sobre BoTorch: experimentos, restricciones de resultado, ensayos pendientes & MIT & Campañas reales con varias personas y máquinas \\
Optuna \texttt{GPSampler}~\cite{optuna} & Optimización bayesiana con PG dentro de Optuna & MIT & Si ya se usa Optuna \\
SMAC3~\cite{smac3} & BO con bosques aleatorios y espacios condicionales & BSD & Variables categóricas y jerárquicas \\
scikit-learn + SciPy & PG y optimización de la adquisición escritos a mano & BSD & Docencia y prototipos (práctica) \\
\bottomrule
\end{tabularx}
\caption{Tecnologías para la optimización bayesiana.}
\label{tab:tec4}
\end{table}

\addtocounter{definicion}{1}%

